% Draft 1, 2026-11-05 (log.md 0067), written from refs-verified.md. Every entry of refs.bib is cited somewhere in the paper.
\section{Related Work}
\label{sec:related}

\paragraph{Exact Counting Under Single Updates.}
Eppstein and Spiro~\cite{eppstein-spiro} maintain the counts of all 3-vertex subgraphs in amortized time proportional to the h-index of the graph, and Eppstein et al.~\cite{eppstein-extended} extend this to 4-vertex subgraphs.
Kara et al.~\cite{kara-tods} maintain a triangle count in $O(\sqrt{m})$ amortized time per update, which is optimal under the OMv conjecture, and Lu and Tao~\cite{lu-tao} give bounds that depend on the arboricity.
Hanauer, Henzinger and Hua~\cite{hhh} count all connected 4-vertex subgraphs with tables that hold only the paths through low-degree vertices, and count through high-degree vertices from the adjacency lists; our mixed tables follow this rule.
Hanauer et al.~\cite{hanauer-dyncount} engineer these approaches on 56 instances and find that, for triangles, scanning the lower-degree endpoint is the fastest choice, while heavy-light tables pay for 4-cycles.
IncGraph~\cite{incgraph} computes the change in all graphlet counts up to five vertices for each edge by local enumeration.
None of these works considers batches, and only IncGraph reaches five vertices.
Hanauer, Henzinger and Schulz~\cite{hanauer-survey} survey the engineering of fully dynamic graph algorithms more broadly.

\paragraph{Batches.}
Dhulipala et al.~\cite{dhulipala-batch-kclique} count triangles and $k$-cliques under batches in parallel, with heavy-light tables, and make each triangle count once by dividing by its number of updated edges.
Liu et al.~\cite{liu-batch-kcore} maintain a low out-degree orientation under batches and bound the work of $k$-clique counting through it.
Makkar et al.~\cite{makkar-dynamic-tc} count triangles on the GPU under batches, by inclusion and exclusion over per-edge intersections, with no work shared between the edges of a batch, and EDTC~\cite{edtc} is a recent GPU counter for dynamic triangles.
Ediger et al.~\cite{ediger-streaming} group a batch by source vertex, but for the update of the data structure rather than for the count.
ParaDyMS~\cite{paradyms} recounts motifs in the modified part of the network after each batch, which is the idea of our recount kernels A4 and A5.

\paragraph{Sharing Inside a Batch.}
Three published methods reuse one row per vertex across the updates of a batch.
DPTL+~\cite{dptl} lists the triangles that a batch adds or removes, in parallel, and builds one bitmap of the neighbourhood of each pivot vertex, which every batch edge at that vertex reuses.
BFCB-IG~\cite{bfcb} counts butterflies, the 4-cycles of a bipartite graph, by building one row of 2-path counts per start vertex over the affected subgraph.
Cheetah~\cite{cheetah} groups each batch into stars for general patterns, caches neighbour lists per star, and reports a mean gain of 1.12 times for triangles and 2.5 to 2.75 times for larger patterns on a cluster of 16 nodes, against Tesseract~\cite{tesseract}, with graph-update time left out.
Gopher~\cite{gopher} shares the set operations within one update instead, and its authors report that this beats Cheetah's sharing across updates.
The idea that this paper measures is therefore published, for triangles, butterflies and general patterns.

\paragraph{Continuous Subgraph Matching.}
General systems for continuous subgraph matching are compared by Sun et al.~\cite{sun-csm-study}, who find that no algorithm wins everywhere and that index updates can dominate the time.
RapidFlow~\cite{rapidflow} reduces each update to batch matching over the affected region, GAMMA~\cite{gamma} matches under batches on the GPU and removes duplicates by a total order on the updates, and Mhedhbi et al.~\cite{mhedhbi-tods} share operators between delta queries.
Ammar et al.~\cite{ammar-delta} state the delta-query rule that reads new versions of earlier relations and old versions of later ones, whose union is exactly the change in the output.
These systems list matches of general patterns, and dedicated counters are faster on triangles~\cite{hanauer-dyncount}.

\paragraph{Static Counting.}
Chiba and Nishizeki~\cite{chiba-nishizeki} list triangles and 4-cycles in time proportional to the arboricity times the number of edges, Alon, Yuster and Zwick~\cite{alon-yuster-zwick} count cycles of a given length, ESCAPE~\cite{escape} counts all 5-vertex subgraphs with a per-vertex row, and Shi, Huang and Shun~\cite{shi-huang-shun} count 5-cycles in $O(m\alpha^2)$ work.
Our recount kernels use the first and the last of these.

\paragraph{Timestamped Streams and This Journal.}
Chanchary and Maheshwari~\cite{chanchary-windowed} answer time-windowed queries on timestamped event graphs, including counts of 2-paths and quadrangles.
Engineering studies of published algorithms, with experiments on real instances, appear in this journal for hypergraph matching~\cite{grossmann-bmatching}, 2-packing sets~\cite{borowitz-2packing} and synchronized planarity~\cite{fink-rutter}, and this paper follows their form.

\begin{table}[t]
  \centering
  \small
  \begin{tabular}{>{\raggedright\arraybackslash}p{2.6cm}>{\raggedright\arraybackslash}p{4.9cm}>{\raggedright\arraybackslash}p{4.9cm}}
    \hline
    Work & What it does & What this paper adds \\
    \hline
    Hanauer et al.~\cite{hanauer-dyncount} & Exact counters for single updates, up to 4 vertices, on 56 instances & Batches, sharing within a batch, and 5-cycles \\
    Dhulipala et al.~\cite{dhulipala-batch-kclique} & Parallel batch counting of cliques with heavy-light tables & Sequential counting of cycles with star rows, against tables, with all costs counted \\
    DPTL+~\cite{dptl} & One bitmap per pivot reused across a batch, for parallel triangle listing & Counting without listing, a degree guard, and up to 5-cycles \\
    BFCB-IG~\cite{bfcb} & One row of 2-path counts per start vertex, for butterflies in bipartite graphs & Odd cycles and general graphs, and a row adjusted for each spoke \\
    Cheetah~\cite{cheetah} & Star grouping with reuse for general patterns, on a cluster, update time left out & A closed-form counting kernel, sequential, against the strongest baselines, with every cost inside \\
    \hline
  \end{tabular}
  \caption{The closest prior work.}
  \label{tab:closest}
\end{table}
